\documentclass[12pt]{article}
\usepackage[utf8]{inputenc}
\usepackage[T1]{fontenc}
\usepackage[margin=1in]{geometry}
\usepackage{setspace}
\usepackage{xcolor}
\usepackage[colorlinks=true, urlcolor=blue!60!black, linkcolor=blue!60!black]{hyperref}
\setlength{\parindent}{2em}
\setlength{\parskip}{1.25em}
\renewcommand{\baselinestretch}{1.0}
\begin{document}
\begin{center}
{\LARGE \textit{Trump v. Cook}, 609 U.S. 528 (2026)\par}
\vspace{0.8em}
{\large Case Brief\par}
\end{center}
\vspace{1em}

\section{Memory Jogger}

In a 5--4 interim-docket decision issued the same day it overruled \textit{Humphrey's Executor}, the Supreme Court upheld the constitutionality of Federal Reserve Governors' statutory ``for cause'' removal protections and kept Governor Lisa Cook in office because the President failed to give her pretermination notice and an opportunity to respond under \href{https://www.law.cornell.edu/supremecourt/text/25A312}{12 U.S.C. \S~242}.

\section{Detailed Case Facts}

The Federal Reserve System, created by the Federal Reserve Act of 1913 and restructured in 1933 and 1935, is a two-tiered central bank comprising 12 privately owned regional Federal Reserve Banks supervised by a seven-member Board of Governors appointed by the President and confirmed by the Senate (\href{https://www.supremecourt.gov/opinions/25pdf/25a312_5468.pdf}{Slip Op. at 5--6}; \href{https://www.courtlistener.com/opinion/10881682/trump-v-cook/}{CourtListener}). Under 12 U.S.C. \S~242, each Governor ``shall hold office for a term of fourteen years from the expiration of the term of his predecessor, unless sooner removed for cause by the President.'' The Federal Reserve also funds its own operations outside annual congressional appropriations (12 U.S.C. \S~243) and sets monetary policy through the Federal Open Market Committee alongside regional bank representatives (12 U.S.C. \S~263).

President Joe Biden appointed Lisa D. Cook to the Board of Governors in 2022 to finish the remaining two years of Janet Yellen's unexpired term, and the Senate confirmed Cook in 2023 to a full 14-year term scheduled to expire in 2038 (\href{https://www.law.cornell.edu/supremecourt/text/25A312}{Cornell LII}). On August 20, 2025, Federal Housing Finance Agency Director William Pulte posted on social media a week-old referral letter to Attorney General Pamela Bondi accusing Cook of mortgage fraud in 2021---prior to her Federal Reserve tenure. Specifically, the letter alleged that on June 18, 2021, Cook obtained a \$203,000 mortgage on a Michigan property declaring it her principal residence for at least one year, and 14 days later obtained a \$540,000 mortgage on a Georgia property also stating it would be her primary residence for one year (\href{https://supreme.justia.com/cases/federal/us/609/25a312/}{Justia}).

Less than 30 minutes after Pulte's post, President Donald J. Trump posted on social media, ``Cook must resign, now!!!'' linking to a news article about the letter, and two days later told reporters he would ``fire her if she doesn't resign'' (\href{https://www.law.cornell.edu/supremecourt/text/25A312}{Cornell LII}). On August 25, 2025, President Trump sent Cook a letter (and posted screenshots online) purporting to remove her ``for cause'' effective immediately, asserting ``reason to believe'' she ``may have made false statements on one or more mortgage agreements'' and declaring that her ``deceitful and potentially criminal conduct'' showed at a minimum ``gross negligence'' incompatible with public confidence in a financial regulator (\href{https://www.supremecourt.gov/opinions/25pdf/25a312_5468.pdf}{Slip Op. at 6--7}; \href{https://www.scotusblog.com/2026/06/court-prevents-trump-from-firing-fed-governor/}{SCOTUSblog}). Cook was the first Governor a President had attempted to fire in the Federal Reserve's 111-year history. Cook denied wrongdoing and received no formal notice of the evidence, no opportunity to respond, and no pretermination hearing; the underlying mortgage documents were never even placed in the judicial record (\href{https://www.brookings.edu/articles/fed-independence-after-trump-v-cook/}{Brookings}).

\section{Procedural History}

\begin{itemize}
    \item \textbf{U.S. District Court for the District of Columbia (No. 1:25-cv-02903):} On August 28, 2025, Cook sued President Trump, Federal Reserve Chair Jerome Powell, and the Board of Governors, seeking declaratory, injunctive, and mandamus relief on the grounds that her firing lacked statutory ``cause'' and violated both statutory and Fifth Amendment Due Process pretermination requirements (\href{https://www.law.cornell.edu/supremecourt/text/25A312}{Cornell LII}). On September 9, 2025, Judge Jia Cobb granted a preliminary injunction barring Powell and the Board from effectuating Cook's removal (\textit{Cook v. Trump}, 804 F. Supp. 3d 14 (D.D.C. 2025)), holding that (1) ``for cause'' under 12 U.S.C. \S~242 applies only to in-office conduct and not pre-appointment mortgages, and (2) Cook held a Fifth Amendment property interest in her office entitling her to pretermination notice and a hearing (\href{https://www.everycrsreport.com/reports/LSB11449.html}{CRS Legal Sidebar LSB11449}).
    \item \textbf{U.S. Court of Appeals for the D.C. Circuit (No. 25-5326):} On September 15, 2025, a divided panel denied the Government's motion for a stay pending appeal (2025 WL 2654786). Judge Garcia, joined by Judge Childs, concurred solely on Cook's Due Process Clause claim. Judge Katsas dissented, arguing that public office is not constitutional ``property,'' that alleged mortgage fraud plainly constitutes ``cause,'' and that Cook failed the demanding standard for nonstatutory \textit{ultra vires} relief (\href{https://supreme.justia.com/cases/federal/us/609/25a312/}{Justia}).
    \item \textbf{Supreme Court of the United States (No. 25A312):} On September 18, 2025, the Government applied for an emergency stay. On October 1, 2025 (606 U.S. 1062), the Court deferred action pending oral argument, which was held on January 21, 2026 (Solicitor General D. John Sauer for applicant; Paul D. Clement for respondent). On June 29, 2026, the Court denied the stay application by a 5--4 vote (\href{https://www.scotusblog.com/2026/06/court-prevents-trump-from-firing-fed-governor/}{SCOTUSblog}).
\end{itemize}

\section{Judicial Votes}

\textbf{Vote:} 5--4 denying the Government's application for a stay (\href{https://www.law.cornell.edu/supremecourt/text/25A312}{Cornell LII}; \href{https://www.courtlistener.com/opinion/10881682/trump-v-cook/}{CourtListener}).

\begin{itemize}
    \item \textbf{Majority (5):} Chief Justice Roberts delivered the opinion of the Court, joined in full by Justices Sotomayor, Kagan, Kavanaugh, and Jackson.
    \item \textbf{Concurrences (2):} Justice Kavanaugh filed a concurring opinion; Justice Jackson filed a concurring opinion.
    \item \textbf{Dissents (4):} Justice Thomas filed a dissenting opinion; Justice Alito filed a dissenting opinion, joined by Justice Gorsuch; Justice Barrett filed a dissenting opinion.
\end{itemize}

\section{Holding}

\textbf{Judgment:} Application for stay denied. The District Court's preliminary injunction keeping Governor Lisa Cook in office pending litigation remains in effect (\href{https://www.supremecourt.gov/opinions/25pdf/25a312_5468.pdf}{Slip Op. at 26--27}).

\textbf{Rule Announced:} Under 12 U.S.C. \S~242, (1) the President's ``for cause'' removal of a Federal Reserve Governor is judicially reviewable; (2) ``cause'' imposes a substantial threshold reflecting the Federal Reserve's unique independence---turning on the seriousness of the alleged misconduct and its nexus to the Governor's official duties, rather than pretext to install a more congenial replacement; (3) federal courts possess equitable authority to issue interim preliminary injunctions protecting an incumbent Governor's possession of office pending merits litigation; (4) the statutory tenure of a fixed term removable only ``for cause'' incorporates the common-law and precedential baseline requiring pretermination notice of the evidence, an avenue to respond (even in writing), and a deadline before a final removal decision; and (5) the Federal Reserve Governors' for-cause removal protection is consistent with Article II of the Constitution given the Nation's distinct historical tradition of independent central banking (\href{https://www.law.cornell.edu/supremecourt/text/25A312}{Cornell LII}; \href{https://www.everycrsreport.com/reports/LSB11449.html}{CRS LSB11449}).

\textbf{Holding vs. Dicta:} The narrow dispositive holding supporting denial of the stay is that the Government is unlikely to succeed on appeal because President Trump failed to afford Cook the statutory pretermination notice and opportunity to respond required by 12 U.S.C. \S~242, rendering his August 25, 2025 removal letter void without needing to reach Cook's Fifth Amendment Due Process claim (\href{https://www.law.cornell.edu/supremecourt/text/25A312}{Cornell LII}). The Court's threshold holdings on judicial reviewability, the availability of interim equitable reinstatement, and the existence of a nonstatutory equitable right of action under \textit{Armstrong v. Exceptional Child Center, Inc.}, 575 U.S. 320 (2015), were necessary steps to deny the stay. By contrast, critics and dissenters (Justices Alito, Gorsuch, and Barrett) argued that the majority's constitutional endorsement of Federal Reserve independence---an issue the Government expressly waived---and its discussion of the substantive ``cause'' standard went beyond what was strictly necessary on an interim stay posture, while footnote 7's caveat that Congress may not necessarily assign the Federal Reserve ``additional regulatory powers that are attenuated from monetary policy'' is classic cautionary dicta (\href{https://www.brookings.edu/articles/fed-independence-after-trump-v-cook/}{Brookings}). Importantly, the Court expressly did \emph{not} decide whether the underlying mortgage allegations, if proven after proper process, would constitute sufficient ``cause'' for removal.

\section{Analysis of Opinions}

\subsection{Majority Opinion (Chief Justice Roberts, joined by Sotomayor, Kagan, Kavanaugh, and Jackson, JJ.)}

Chief Justice Roberts framed the dispute through a sweeping historical narrative tracing American central-bank independence from the Revolutionary-era Bank of North America through the First Bank of the United States (1791) and Second Bank of the United States (1816)---both largely insulated from presidential control---to the financial panics following Andrew Jackson's veto of the Second Bank's charter and the eventual creation of the Federal Reserve in 1913 (\href{https://www.law.cornell.edu/supremecourt/text/25A312}{Cornell LII}). Applying the stay standard of \textit{Nken v. Holder}, 556 U.S. 418 (2009), and \textit{Hollingsworth v. Perry}, 558 U.S. 183 (2010), the Court stopped at the first factor---likelihood of success on the merits---and rejected all three of the Government's threshold arguments.

First, the Court held that ``for cause'' determinations under 12 U.S.C. \S~242 are judicially reviewable. Citing \textit{Loper Bright Enterprises v. Raimondo}, 603 U.S. 369 (2024), and common-law authorities back to \textit{Bagg's Case} (1615), Roberts explained that delegating removal authority to the President does not commit the statutory boundary of ``cause'' to his unreviewable discretion, distinguishing \textit{Reagan v. United States}, 182 U.S. 419 (1901), where Congress had failed to prescribe any causes at all (\href{https://www.law.cornell.edu/supremecourt/text/25A312}{Cornell LII}).

Second, the Court rejected both extremes on the meaning of ``cause.'' The Government's low bar (any concern about conduct, fitness, or integrity) would invite pretextual firings destructive of both the fact and appearance of Fed independence, whereas Cook's rigid bar (strictly in-office inefficiency, neglect, or malfeasance) lacked textual grounding in \S~242. Drawing on common-law principles, Roberts held that ``cause'' sets a ``substantial threshold'' turning on the seriousness of the alleged misconduct and its nexus to official duties, asking whether the stated reason genuinely implies unfitness or is merely a pretext to install a ``more congenial'' replacement (\href{https://www.law.cornell.edu/supremecourt/text/25A312}{Cornell LII}).

Third, the Court upheld the District Court's authority to enter an interim preliminary injunction keeping Cook in office. Distinguishing \textit{In re Sawyer}, 124 U.S. 200 (1888), and \textit{White v. Berry}, 171 U.S. 366 (1898), Roberts relied on 19th-century equity treatises (High and Pomeroy) and modern scholarship to distinguish final adjudication of title to office (historically reserved to law via \textit{quo warranto} or mandamus) from interim equitable protection of \textit{de facto} incumbents pending litigation. In footnote 3, Roberts likewise rejected Justice Thomas's lack-of-cause-of-action argument by invoking traditional \textit{Ex parte Young}-style equitable suits to enjoin unlawful federal official action (\href{https://www.law.cornell.edu/supremecourt/text/25A312}{Cornell LII}).

Turning to the narrow ground of decision, Roberts held that Congress enacted \S~242 in 1913 (and reenacted it in 1935) against the settled rule of \textit{Reagan} and \textit{Shurtleff v. United States}, 189 U.S. 311 (1903), that where an office is for a fixed term removable only for cause, ``notice and hearing are essential.'' While a Governor is not entitled to an oral hearing or judicial trial, the President must provide ``some explanation of the evidence at issue, some avenue for a response, and a deadline by which a response would be due'' before deciding (\href{https://www.supremecourt.gov/opinions/25pdf/25a312_5468.pdf}{Slip Op. at 24--25}). A social-media post demanding Cook's resignation fell far short. Finally, responding to Justice Thomas's constitutional attack, the majority explicitly upheld the constitutionality of Federal Reserve Governors' for-cause protection under Article II, concluding that---like the First and Second Banks---the Fed maintains the founding-era balance shielding monetary policy from political manipulation without being ``trapped in amber'' (\href{https://www.law.cornell.edu/supremecourt/text/25A312}{Cornell LII}).

\subsection{Concurring Opinion (Justice Kavanaugh)}

Justice Kavanaugh joined the Court's opinion in full to emphasize two points. First, the interim ruling does not decide whether President Trump may ultimately remove Cook for cause once the facts are developed through proper procedure. Second, Kavanaugh vigorously defended the majority's decision to settle the Federal Reserve's constitutional status immediately after \textit{Trump v. Slaughter}. Emphasizing that the Government itself conceded the Fed's unique historical status at oral argument, Kavanaugh warned that leaving the question open would inject dangerous uncertainty into U.S. and global markets and invite immediate political turmoil (\href{https://www.law.cornell.edu/supremecourt/text/25A312}{Cornell LII}; \href{https://www.paulweiss.com/insights/client-memos/supreme-court-overrules-humphrey-s-executor-but-preserves-the-fed-s-independence}{Paul, Weiss}).

\subsection{Concurring Opinion (Justice Jackson)}

Justice Jackson joined the majority opinion but wrote separately to argue that the stay application failed overwhelmingly on the equitable factors alone. She rejected the Government's claim that merely keeping an official in office over the President's objection constitutes irreparable injury, calling it ``just another species of the far-fetched contention that the President [is] injured whenever he is prevented from doing as he wishes'' (\href{https://www.law.cornell.edu/supremecourt/text/25A312}{Cornell LII}). Weighing the public interest, Jackson stressed the bipartisan consensus of economists and former Treasury Secretaries that even the perception of political intimidation at the Fed risks inflation and financial panic, making the equities ``not a close call.''

\subsection{Dissenting Opinion (Justice Thomas)}

Justice Thomas dissented across the board. First, he argued that Cook had no Fifth Amendment ``property'' right in a public office (\textit{Taylor v. Beckham}, 178 U.S. 548 (1900)). Second, he argued that \S~242's bare ``for cause'' text sets the weakest possible statutory standard, easily satisfied by apparent federal mortgage fraud. Third, he attacked the majority's statutory notice-and-hearing requirement as an atextual invention resting on a ``single line of dicta'' in \textit{Reagan} and \textit{Shurtleff}, pointing to dozens of statutes where Congress expressly required notice and a hearing when it wanted them (\href{https://www.law.cornell.edu/supremecourt/text/25A312}{Cornell LII}). Fourth, breaking with the Government's concession, Thomas declared the Federal Reserve Act's removal restriction flatly unconstitutional under Article II and \textit{Trump v. Slaughter}: unlike the private First and Second Banks, which exercised no sovereign power, the Federal Reserve Board is a modern executive agency modeled on the German Reichsbank that issues binding rules, levies \$1 million daily fines, and bans individuals from banking. Finally, Thomas argued that federal courts lacked a statutory cause of action, a waiver of sovereign immunity, and equitable jurisdiction (\textit{Sawyer}, \textit{White}, \textit{Sampson v. Murray}) to enjoin the President's ``conclusive and preclusive'' removal power (\href{https://supreme.justia.com/cases/federal/us/609/25a312/}{Justia}).

\subsection{Dissenting Opinion (Justice Alito, joined by Justice Gorsuch)}

Justice Alito, joined by Justice Gorsuch, dissented primarily on procedural and institutional-restraint grounds. Criticizing the Court for freezing lower-court litigation for nine months and issuing a sweeping merits-style opinion on an undeveloped 21-day record that lacked even the mortgage documents, Alito argued the Court should have acted solely on the two grounds decided below: (1) the District Court erred in limiting ``for cause'' to in-office conduct, and (2) both lower courts erred in finding a Fifth Amendment property interest in a public office (\href{https://www.law.cornell.edu/supremecourt/text/25A312}{Cornell LII}; \href{https://www.everycrsreport.com/reports/LSB11449.html}{CRS LSB11449}). Because the lower courts' actual rationales were wrong, Alito would have granted the stay and left all novel statutory and constitutional questions to percolate below.

\subsection{Dissenting Opinion (Justice Barrett)}

Justice Barrett echoed Alito's objection to deciding a host of unbriefed and waived questions on the interim docket, singling out the majority's two-paragraph constitutional validation of the Federal Reserve as ``entirely outside the scope of this case'' and ``in serious tension with \textit{Trump v. Slaughter}'' (\href{https://www.law.cornell.edu/supremecourt/text/25A312}{Cornell LII}). She questioned how history could simultaneously justify \textit{Slaughter}'s categorical rule and \textit{Cook}'s carveout, and whether the Fed's non-monetary regulatory powers fit within that exception. Barrett also identified a remedial mismatch: because the District Court's injunction barred removal ``on the basis of the grounds stated in the President's letter of August 25, 2025'' (mortgage fraud), denying the stay left in place an overbroad injunction preventing the President from trying again even with full process, notwithstanding the majority's narrow gloss in footnote 9.

\section{Examples: Future Applications}

\begin{enumerate}
    \item \textbf{Same-Side (Summary Firing Over Interest-Rate Votes):} Suppose a President publicly criticizes a Federal Reserve Governor for refusing to cut interest rates before an election and then fires the Governor the next day via executive order citing ``gross economic mismanagement,'' without providing evidence or a chance to respond. Under \textit{Cook}, a federal court would preliminarily enjoin the removal: the President violated \S~242's mandatory pretermination notice-and-response process, and ``mismanagement'' rooted in monetary-policy disagreement fails the ``substantial threshold'' for cause and is a textbook pretext to obtain a ``more congenial'' replacement.
    \item \textbf{Same-Side (Pretextual Firing Over Minor Technical Filing Error):} Suppose a President gives a Federal Reserve Governor written notice and 10 days to respond to an allegation that, six years before joining the Board, the Governor filed a routine municipal property-tax disclosure form two weeks late and paid a \$150 administrative late fee, and then fires the Governor ``for cause.'' Even though procedural notice was provided, a reviewing court under \textit{Cook} would likely invalidate the removal because an isolated, minor, pre-office administrative oversight lacks both the seriousness and the nexus to official duties required to demonstrate genuine unfitness for the Board.
    \item \textbf{Opposite-Side (Removal After Felony Conviction for Insider Trading Following Notice and Opportunity to Respond):} Suppose a sitting Federal Reserve Governor is indicted and pleads guilty to leaking nonpublic Federal Open Market Committee rate decisions to a hedge fund, and the President provides the Governor with written notice of the plea and investigative file, gives 14 days to submit a written defense, reviews the submission, and issues a written removal order. Under \textit{Cook}, this removal would be upheld: the President satisfied \S~242's minimal written process, and serious criminal misconduct directly implicating the integrity of monetary policy plainly satisfies the ``substantial threshold'' for cause.
    \item \textbf{Opposite-Side (At-Will Removal of an SEC or FDIC Commissioner):} Suppose a President fires a Commissioner of the Securities and Exchange Commission or a Director of the Federal Deposit Insurance Corporation solely for policy disagreements, without notice or cause. Unlike a Federal Reserve Governor, those officials lack the Fed's unique historical lineage tied to the First and Second Banks of the United States and monetary policy; under the companion decision in \textit{Trump v. Slaughter}, their removal protections are unconstitutional under Article II, so \textit{Cook}'s protections and interim reinstatement would not apply.
    \item \textbf{Fence-Sitter (Removal of a Fed Governor After Written Process for Disputed Pre-Office Financial Misconduct):} Suppose the White House restarts proceedings against a Federal Reserve Governor by sending formal written notice of pre-office mortgage discrepancies along with loan documents, affording 20 days for written rebuttal; the Governor responds with closing disclosures showing the second home was declared a vacation property on other forms and argues any error was inadvertent, yet the President still finds ``gross negligence'' and removes her. This sits squarely on the fence left open by \textit{Cook}: the President satisfied \S~242's procedural floor, forcing a court to decide the unresolved substantive questions of how much deference the President's factual findings receive, whether pre-office negligence (short of intentional fraud) clears the ``substantial threshold,'' and whether an interim injunction should issue once process has been afforded (\href{https://www.brookings.edu/articles/fed-independence-after-trump-v-cook/}{Brookings}).
\end{enumerate}

\section{Critique}

\textbf{Originalist and Textualist Critiques:} Strict unitary-executive scholars and Justices Thomas and Barrett criticized \textit{Cook} as methodologically incoherent when read beside \textit{Trump v. Slaughter}. From an originalist standpoint, Thomas argued that the First and Second Banks of the United States were private corporations exercising no sovereign regulatory power, making the majority's historical analogy to the modern Federal Reserve Board---an administrative agency that promulgates binding banking rules and levies civil penalties---an ahistorical fiction (\href{https://www.law.cornell.edu/supremecourt/text/25A312}{Cornell LII}). Textualists likewise faulted Roberts for reading a notice-and-hearing requirement into 12 U.S.C. \S~242 despite Congress's omission of procedural language that it explicitly included in contemporaneous statutes like the 1935 National Labor Relations Act. Former Federal Reserve Governor and Harvard Law Professor Daniel K. Tarullo similarly observed that ``[h]istorians will be struck by the superficial nature of the Chief Justice's account'' and that the Fed carveout reflects a pragmatic policy judgment rather than rigorous originalism (\href{https://www.brookings.edu/articles/fed-independence-after-trump-v-cook/}{Brookings}). Procedural formalists (Alito, Gorsuch, and Barrett) also criticized the Court for deciding a waived constitutional issue on the emergency docket without full lower-court percolation.

\textbf{Progressive and Doctrinal Critiques:} While progressive commentators and economists welcomed the preservation of Federal Reserve independence and the Court's willingness to police pretextual firings, they highlighted deep doctrinal tensions with \textit{Slaughter}. If the Court recognizes that insulating monetary policy from short-term presidential interference serves vital public stability and that ``[n]ot only the fact of independence but also the appearance of independence'' matters, critics ask why the same functionalist logic does not protect agencies overseeing nuclear safety, consumer markets, or financial regulation (\href{https://www.brookings.edu/articles/fed-independence-after-trump-v-cook/}{Brookings}; \href{https://www.paulweiss.com/insights/client-memos/supreme-court-overrules-humphrey-s-executor-but-preserves-the-fed-s-independence}{Paul, Weiss}). Moreover, Tarullo and the Congressional Research Service noted significant unresolved vulnerabilities: the majority's minimalist procedural rule (written notice and response) leaves unclear what factual record courts will review after a renewed removal effort, and footnote 7's warning about ``additional regulatory powers that are attenuated from monetary policy'' invites future challenges to the Fed's bank-supervision and consumer-fee authorities (\href{https://www.everycrsreport.com/reports/LSB11449.html}{CRS LSB11449}).

\textbf{Underlying Competing Values:} \textit{Cook} exposes a sharp collision between democratic accountability under a unitary executive and institutional insulation needed for credible, long-horizon monetary policy. The majority and Justice Kavanaugh prioritized macroeconomic stability, market confidence, and historical pragmatism---refusing to risk global financial turmoil in service of doctrinal purity---whereas Justice Thomas prioritized constitutional formalism and electoral accountability, and Justices Alito, Gorsuch, and Barrett prioritized judicial minimalism and disciplined adherence to interim-docket boundaries.

\section{Key Quotations}

\begin{quote}
``To accept any one of those arguments would in effect transform the Federal Reserve's for-cause protection into at-will employment---an interpretive leap out of step with the statute Congress enacted and our Nation's tradition of central banking protected from political interference.''
\end{quote}
\noindent --- Chief Justice Roberts, Opinion of the Court (\href{https://www.law.cornell.edu/supremecourt/text/25A312}{Cornell LII}; \href{https://www.supremecourt.gov/opinions/25pdf/25a312_5468.pdf}{Slip Op. at 9}).

\begin{quote}
``Not only the fact of independence but also the appearance of independence is key to the Federal Reserve's design. \ldots{} Without such constraints in place, any perceived or alleged misstep (past or present) could provide a ready pretext for a Governor's removal---a fact that he would surely know, and that would surely weigh on him as he decided what to say and how to vote. Nothing could be more corrosive of the independence that Congress sought to preserve.''
\end{quote}
\noindent --- Chief Justice Roberts, Opinion of the Court (\href{https://www.law.cornell.edu/supremecourt/text/25A312}{Cornell LII}; \href{https://www.supremecourt.gov/opinions/25pdf/25a312_5468.pdf}{Slip Op. at 14--15}).

\begin{quote}
``At minimum, Cook was entitled to some explanation of the evidence at issue, some avenue for a response, and a deadline by which a response would be due.''
\end{quote}
\noindent --- Chief Justice Roberts, Opinion of the Court (\href{https://www.law.cornell.edu/supremecourt/text/25A312}{Cornell LII}; \href{https://www.supremecourt.gov/opinions/25pdf/25a312_5468.pdf}{Slip Op. at 24}).

\begin{quote}
``Even temporary uncertainty about the status of the Federal Reserve could spark political upheaval, including confusion about whether the President could immediately remove multiple Governors at will, as well as turmoil in the U. S. and world economies.''
\end{quote}
\noindent --- Justice Kavanaugh, Concurring Opinion (\href{https://www.law.cornell.edu/supremecourt/text/25A312}{Cornell LII}; \href{https://supreme.justia.com/cases/federal/us/609/25a312/}{Justia}).

\begin{quote}
``Although the Court expresses concern that the President removed a Board member for `the first time in the Federal Reserve's 111-year history,' it expresses no such concern that it today upholds an injunction against the President's removal of an executive officer for the first time in the Constitution's 237-year history.''
\end{quote}
\noindent --- Justice Thomas, Dissenting Opinion (\href{https://www.law.cornell.edu/supremecourt/text/25A312}{Cornell LII}; \href{https://supreme.justia.com/cases/federal/us/609/25a312/}{Justia}).


\end{document}
